% from main.tex
%%%%%%%%%%  scan idx 171  |  folio 158  %%%%%%%%%%
trace se prolonge de manière unique en un \(S\)-foncteur de \(\mathcal E(\Dpx(S)_{\parf})^{\mathrm{is}}\) dans \(\mathcal A\).

Pour \(E\in\Ob\Parf(U)\) (\(U\in\Ob S\)), on note encore
\[
  \Tr_E:\Hom(E,E)\longrightarrow H^0(U,A)
\]
l'homomorphisme correspondant, qui est un homomorphisme de \(H^0(U,A)\)-modules. On voit de même que la définition du cup-produit (8.2), ainsi que les assertions (8.3) et (8.4) se généralisent au cas de complexes parfaits (non nécessairement de tor-dimension finie).

\textbf{Proposition 8.6.} Soient \(U\in\Ob S\), \(E,F\in\Parf(U)\), \(u\in\Hom(E,E)\), \(v\in\Hom(F,F)\), d'où \(u\otL v\in\Hom(E\otL F,E\otL F)\). On a
\[
  \Tr_{E\otL F}(u\otL v)=\Tr_E(u)\,\Tr_F(v).
\]

\emph{Preuve.} La question étant locale, on peut supposer grâce à (8.4) et (4.10), que \(u\) et \(v\) sont des flèches de \(K^b(\mathcal C_0)\). Utilisant (8.1.2), on se ramène au cas où \(u\) et \(v\) sont des flèches de \(\mathcal C_0\). On termine grâce à (8.1.1).

\section*{Bibliographie}
\begin{description}[leftmargin=2.5cm,style=nextline]
\item[{[H]}]
